% ----------------------------------------------------------
% Introdução
% ----------------------------------------------------------
\chapter{Introdução}

A Indústria 4.0 impulsiona a demanda por robôs móveis autônomos (AMR) em ambientes logísticos e industriais \cite{Anand_2024_intro, IFR2025, braun2022robot}. No entanto, plataformas comerciais frequentemente operam como ``caixas-pretas'', limitando o aprendizado prático de controle embarcado, cinemática e percepção sensorial \cite{patel2023open, jolles2022open}. O \textit{Robot Operating System 2} (ROS 2) se estabelece como \textit{framework} aberto e modular para superar essa barreira \cite{macenski2022robot}. No âmbito do Instituto Federal do Piauí (IFPI), \citeonline{vilanova2022aeronave} e \citeonline{Dantas_2022} desenvolveram plataformas robóticas educacionais baseadas em ROS 2, sendo o Roomba® 614 a base de hardware adotada pelo segundo.

Este artigo descreve a evolução dessa plataforma: um robô seguidor de linha autônomo adaptado do iRobot Roomba® 614, com controlador PID discreto e \textit{middleware} ROS 2 Humble Hawksbill executados diretamente em um Raspberry Pi 4 Model B via \textit{Edge Computing} (Computação de Borda) \cite{vega2025bridging, kortelainen2023short}. A centralização do processamento na placa elimina microcontroladores intermediários, reduzindo latência e custo, e cria uma plataforma didática aberta para aplicações em logística e ensino de engenharia.